% Generated by analysis/analysis.py -- do not edit by hand.
\begin{table}[H]
  \centering
  \renewcommand{\arraystretch}{1.3}
  \begin{tabular}{|c|c|c|c|}
    \hline
    גודל & מספר מדידות & ערך & שגיאה יחסית \\ \hline
    $m$ & 4 & $1.2700 \pm 0.0029\U{g}$ & $0.23\%$ \\ \hline
    $t$ & 10 & $18.623 \pm 0.041\U{s}$ & $0.22\%$ \\ \hline
    $d$ & 4 & $5.938 \pm 0.028\U{mm}$ & $0.47\%$ \\ \hline
    $D_{ch}$ & 4 & $6.025 \pm 0.014\U{mm}$ & $0.22\%$ \\ \hline
    $D$ & 4 & $19.955 \pm 0.017\U{mm}$ & $0.084\%$ \\ \hline
    $L_0$ & 1 & $20.000 \pm 0.029\U{cm}$ & $0.14\%$ \\ \hline
    $a$ & מחושב & $3.0125 \pm 0.0068\U{mm}$ & $0.22\%$ \\ \hline
    $r_m$ & מחושב & $2.969 \pm 0.014\U{mm}$ & $0.47\%$ \\ \hline
    $w$ & מחושב & $6.965 \pm 0.011\U{mm}$ & $0.16\%$ \\ \hline
  \end{tabular}
  \caption{הגדלים ששימשו לחישוב המוליכות החשמלית לפי נוסחה~\eqref{eq:sigma}, ושגיאותיהם. $a$, $r_m$ ו-$w$ חושבו מהקטרים שנמדדו. השגיאה ב-$\sigma$ חושבה ישירות מ-$\Delta D_{ch}$ ומ-$\Delta D$, לפי נוסחה~\eqref{eq:dsigma}, ולא מ-$\Delta a$ ומ-$\Delta w$.}
  \label{tab:a-values}
\end{table}
